\documentclass{beamer}

% --- для схем (TikZ) ---
\usepackage{tikz}
\usetikzlibrary{arrows.meta,positioning,shapes.geometric,fit}

% --- для графиков/таблиц (PGFPlots) ---
\usepackage{pgfplots}
\pgfplotsset{compat=1.18}

% --- аккуратные таблицы ---
\usepackage{booktabs}

\usetheme{Madrid}
\usecolortheme{default}

\usepackage[utf8]{inputenc}
\usepackage[english,russian]{babel}

\title{Название презентации}
\author{Автор}
\date{\today}

\begin{document}

\frame{\titlepage}

\begin{frame}
\frametitle{Содержание}
\tableofcontents
\end{frame}

\section{Первый раздел}


\begin{figure}{}
\centering
\input{figures/figure1}
\caption{An example of using the proof generation model to generate a proof.}
\end{figure}

\begin{frame}
\frametitle{Первый слайд}
\begin{itemize}
\item Пункт 1
\item Пункт 2
\item Пункт 3
\end{itemize}
\end{frame}

\section{Второй раздел}

\begin{frame}
\frametitle{Второй слайд}
Содержание слайда
\end{frame}

% =================================================
%                Денис
% =================================================
\begin{frame}{Что такое LLM и почему это важно}
\textbf{LLM} — большая языковая модель (Transformer), обученная на больших корпусах текста/кода:
\begin{itemize}
  \item генерирует текст и код;
  \item умеет «достраивать» последовательности;
  \item может быть дообучена на специализированных данных (fine-tuning).
\end{itemize}

\vspace{2mm}
В статье используется \textbf{Minerva} — модель, предварительно обученная на математических текстах.
\end{frame}

\begin{frame}{Whole-proof generation: идея}
Вместо пошагового поиска:
\begin{itemize}
  \item вход: формулировка теоремы;
  \item выход: \textbf{полный proof script}.
\end{itemize}

\vspace{2mm}
Дальше proof assistant просто проверяет:
\begin{itemize}
  \item нет ошибок $\Rightarrow$ теорема доказана;
  \item есть ошибки $\Rightarrow$ пробуем другой сэмпл или включаем repair.
\end{itemize}
\end{frame}

\end{document}